\documentclass[10pt,a4paper,landscape]{article}

\usepackage[a4paper, landscape, top=0.8cm, bottom=0.8cm, left=0.6cm, right=0.6cm]{geometry}
\usepackage{multicol}
\usepackage{amsmath, amssymb}
\usepackage{graphicx}
\usepackage{booktabs}
\usepackage{xcolor}
\usepackage{mdframed}
\usepackage{enumitem}
\usepackage{microtype}
\usepackage[T1]{fontenc}
\usepackage[utf8]{inputenc}

\setlength{\columnseprule}{0.4pt}
\setlength{\columnsep}{0.8cm}
\setlength{\parindent}{0pt}
\setlength{\parskip}{2pt}

\setlist[itemize]{leftmargin=*, itemsep=0pt, topsep=1pt}
\setlist[enumerate]{leftmargin=*, itemsep=0pt, topsep=1pt}

\definecolor{headbg}{RGB}{30,80,140}
\definecolor{exbg}{RGB}{220,235,255}
\definecolor{alertbg}{RGB}{255,230,220}

\newcommand{\secbox}[1]{%
  \vspace{3pt}%
  \noindent\colorbox{headbg}{\makebox[\dimexpr\linewidth-2\fboxsep\relax][l]{\color{white}\textbf{\small #1}}}%
  \vspace{2pt}%
}

\newcommand{\formula}[1]{%
  \begin{mdframed}[backgroundcolor=exbg, linewidth=0pt, innertopmargin=3pt, innerbottommargin=3pt, innerleftmargin=4pt, innerrightmargin=4pt]
  \small #1
  \end{mdframed}
}

\newcommand{\alerta}[1]{%
  \begin{mdframed}[backgroundcolor=alertbg, linewidth=0pt, innertopmargin=3pt, innerbottommargin=3pt, innerleftmargin=4pt, innerrightmargin=4pt]
  \small #1
  \end{mdframed}
}

\begin{document}
\footnotesize

\begin{center}
  {\large\textbf{Cheat Sheet: DLP, Diffie-Hellman y ECC}} \quad
  {\small Matemática Discreta --- UNI 2026}\\[2pt]
  {\footnotesize Eliaz S. Bobadilla Camarena \quad Raúl Papuico Palomino \quad Carlo S. Bustamante Martínez \quad Carlos A. Quispe Luna}
\end{center}
\vspace{2pt}
\hrule
\vspace{4pt}

\begin{multicols}{3}

%% =============================================
%% COLUMNA 1: FUNDAMENTOS
%% =============================================

\secbox{1. Grupo $\mathbb{Z}_p^*$ y Raíz Primitiva}

\textbf{Grupo $\mathbb{Z}_p^*$:} conjunto $\{1, 2, \dots, p{-}1\}$ con multiplicación módulo $p$ (primo).

\textbf{Orden de $g$:} menor $k>0$ tal que $g^k \equiv 1 \pmod{p}$.

\formula{
  $g$ es \textbf{raíz primitiva} mod $p$ si\\
  $\text{ord}(g) = \phi(p) = p-1$\\
  $\Rightarrow \{g^1, g^2, \dots, g^{p-1}\} = \mathbb{Z}_p^*$
}

\textbf{Ejemplo} ($p=11$, $g=2$):

\begin{center}
\begin{tabular}{c|cccccccccc}
\small$x$ & 1&2&3&4&5&6&7&8&9&10\\\hline
\small$2^x$ & 2&4&8&5&10&9&7&3&6&1
\end{tabular}
\end{center}
Genera los 10 elementos $\Rightarrow$ $g=2$ es raíz primitiva.

\vspace{3pt}
\secbox{2. Problema del Logaritmo Discreto (DLP)}

\formula{
  Dado $h \equiv g^x \pmod{p}$,\\
  hallar $x \in \{0,\dots,p{-}2\}$
}

\begin{itemize}
  \item \textbf{Fácil:} calcular $g^x \pmod{p}$ (exponenciación binaria, $O(\log x)$ multiplicaciones)
  \item \textbf{Difícil:} invertir — no hay algoritmo eficiente para $p$ grande
\end{itemize}

\textbf{Exponenciación binaria rápida:}
\begin{enumerate}
  \item Escribir $x$ en binario
  \item Aplicar \textit{doblar y multiplicar}
\end{enumerate}
\textit{Ej:} $2^{10} \pmod{11}$: $10=(1010)_2$
$\rightarrow 2^2=4,\; 4^2=5,\; 5\cdot 2^0=5,\; 5^2=3,\; 3\cdot 2^0=3$... (resultado: $1$)

\vspace{3pt}
\secbox{Propiedad clave de índices}

\formula{
  $\text{ind}_g(c^t) \equiv t \cdot \text{ind}_g(c) \pmod{p{-}1}$
}

Esto hace que $(g^a)^b = (g^b)^a = g^{ab}$ — base del intercambio de claves.

%% =============================================
%% COLUMNA 2: DIFFIE-HELLMAN CLÁSICO
%% =============================================

\secbox{3. Protocolo Diffie-Hellman (DH)}

\textbf{Setup público:} primo $p$, raíz primitiva $g$. Ambos conocen $(p,g)$.

\vspace{2pt}
\begin{enumerate}
  \item \textbf{Alice} elige secreto $a$, publica $A = g^a \pmod{p}$
  \item \textbf{Bob} elige secreto $b$, publica $B = g^b \pmod{p}$
  \item \textbf{Alice} calcula $s = B^a \pmod{p}$
  \item \textbf{Bob} calcula $s = A^b \pmod{p}$
\end{enumerate}

\formula{
  $s = (g^b)^a = (g^a)^b = g^{ab} \pmod{p}$\\
  \textbf{Ambos obtienen el mismo secreto} $s$
}

\textbf{Ejemplo numérico} ($p=11,\; g=2$):

\begin{center}
\begin{tabular}{ll}
\toprule
Alice ($a=3$) & Bob ($b=4$) \\\midrule
$A = 2^3 = \mathbf{8}$ & $B = 2^4 = \mathbf{5}$ \\
$s = 5^3 = 125 \equiv \mathbf{4}$ & $s = 8^4 = 4096 \equiv \mathbf{4}$ \\\bottomrule
\end{tabular}
\end{center}
Clave compartida: $s = \mathbf{4}$ \checkmark

\vspace{3pt}
\secbox{4. Problemas CDH y DDH}

\textbf{CDH} (Computational DH):
\begin{itemize}
  \item Conoces: $g,\; p,\; g^a,\; g^b$
  \item Hallar: $g^{ab}$ \quad $\leftarrow$ inviable
\end{itemize}

\alerta{
  \textbf{Trampa común:} $g^a \cdot g^b = g^{a+b} \neq g^{ab}$
}

\textbf{DDH} (Decisional DH): distinguir $g^{ab}$ de valor aleatorio en $\mathbb{Z}_p^*$. Si no se puede $\Rightarrow$ \textbf{seguridad semántica}.

\vspace{3pt}
\secbox{5. Ataque Man-in-the-Middle (MitM)}

\begin{center}
Alice $\leftrightarrow$ \textbf{Eva} $\leftrightarrow$ Bob
\end{center}
Eva intercepta $A$ y $B$, crea dos secretos separados. DH solo es seguro si los valores públicos están \textbf{autenticados} (ej. firmas RSA, certificados TLS).

%% =============================================
%% COLUMNA 3: ECC Y COMPARACIÓN
%% =============================================

\secbox{6. Curvas Elípticas (ECC)}

\formula{
  $y^2 \equiv x^3 + ax + b \pmod{p}$\\
  Condición: $4a^3 + 27b^2 \not\equiv 0 \pmod{p}$
}

\begin{itemize}
  \item Punto al infinito $\mathcal{O}$ = elemento neutro
  \item \textbf{Suma de puntos:} operación geométrica (reflexión)
  \item \textbf{Multiplicación escalar:} $kG = \underbrace{G+G+\cdots+G}_{k}$
\end{itemize}

\secbox{7. ECDLP}

\formula{
  Dado $Q = kG$, hallar $k$ — inviable
}

Sin algoritmos sub-exponenciales conocidos para curvas bien elegidas $\Rightarrow$ más duro que DLP clásico.

\vspace{3pt}
\secbox{8. Protocolo ECDH}

\begin{enumerate}
  \item Setup: curva $E$, punto base $G$, primo $p$
  \item Alice: secreto $a$, publica $A = aG$
  \item Bob: secreto $b$, publica $B = bG$
  \item Secreto: $S = aB = bA = abG$
\end{enumerate}
Usar coordenada $x$ de $S$ + KDF (ej. SHA-256) $\rightarrow$ clave AES.

\vspace{3pt}
\secbox{9. DH Clásico vs. ECDH}

\begin{center}
\begin{tabular}{lcc}
\toprule
\textbf{Seguridad} & \textbf{DH} & \textbf{ECDH} \\\midrule
80 bits  & 1024 b & 160 b \\
128 bits & 3072 b & \textbf{256 b} \\
256 bits & 15360 b & 512 b \\\bottomrule
\end{tabular}
\end{center}

ECDH = misma seguridad, claves $\sim$10$\times$ más cortas.

\vspace{3pt}
\secbox{10. Uso real}

\begin{itemize}
  \item \textbf{TLS 1.3} (HTTPS): ECDH para intercambio de clave de sesión
  \item \textbf{Signal / WhatsApp}: X25519 (curva elíptica)
  \item \textbf{SSH}: DH o ECDH en handshake
  \item \textbf{PGP}: RSA o ECC para cifrado de clave
\end{itemize}

\alerta{
  \textbf{Flujo completo:}\\
  $p,g$ públicos $\to$ claves $A,B$ $\to$ secreto $g^{ab}$ $\to$ KDF $\to$ AES/Vernam
}

\end{multicols}

\end{document}
